\chapter{TypeScript}

JavaScript ha un sistema di tipi peculiare, che abbiamo imparato ad
\sout{amare} accettare.

È \term{tipizzato dinamicamente}:

\begin{affianco}[0.5][c]
\begin{lstlisting}[style=js]
let x;
console.log(typeof(x));   // undefined
x = 1.25e7;
console.log(typeof(x));   // number
x = "Hello Web Technologies!";
console.log(typeof(x));   // string
\end{lstlisting}
\risultato
\begin{immagine}[La stessa variabile cambia tipo nel tempo.]
\begin{tikzpicture}[
    t/.style={rounded corners=2pt, draw=primary!60!black, fill=primary!12!white,
      line width=0.7pt, font=\FiraCodeNL\scriptsize, minimum width=1.7cm, minimum height=0.5cm},
    a/.style={-{Stealth[length=4pt,width=3.5pt]}, line width=0.7pt, draw=neutral!45!black}]
  \node[t] (a) at (0,0) {undefined};
  \node[t] (b) at (2.3,0) {number};
  \node[t] (c) at (4.6,0) {string};
  \draw[a] (a) -- (b); \draw[a] (b) -- (c);
  \node[mnota, anchor=north] at (a.south) {\code{let x}};
  \node[mnota, anchor=north] at (b.south) {\code{x = 1.25e7}};
  \node[mnota, anchor=north] at (c.south) {\code{x = "Hello\ldots"}};
\end{tikzpicture}
\end{immagine}
\end{affianco}

Ed è anche \term{tipizzato debolmente}, con una spiccata predilezione per le
conversioni implicite:

\begin{affianco}[0.6][c]
\begin{lstlisting}[style=js]
let y = "100" - 1 + "42";
// string - number -> number, poi number + string -> string
console.log(`${y}: ${typeof(y)}`);   // 9942: string
\end{lstlisting}
\risultato
\begin{immagine}[Due conversioni implicite in una riga.]
\begin{tikzpicture}[
    v/.style={mcod},
    a/.style={-{Stealth[length=4pt,width=3.5pt]}, line width=0.7pt, draw=neutral!45!black},
    l/.style={mnota, above, text=accent!70!black}]
  \node[v] (a) at (0,0) {"100" - 1};
  \node[v] (b) at (2.2,0) {99 + "42"};
  \node[v] (c) at (4.3,0) {"9942"};
  \draw[a] (a) -- node[l]{\(-\) numerico} (b);
  \draw[a] (b) -- node[l]{\(+\) concatena} (c);
\end{tikzpicture}
\end{immagine}
\end{affianco}

\section{Il problema}

Queste caratteristiche rendono il linguaggio pratico per piccole manipolazioni
di pagine web: funziona e basta, senza bisogno di essere prolissi dichiarando
tipi o conversioni esplicite.

Man mano che la complessità dei programmi cresce, però:

\begin{itemize}
  \item è più facile introdurre bug insidiosi;
  \item l'esperienza di sviluppo peggiora: poco supporto dall'IDE, pochissimi
        errori individuati prima dell'esecuzione;
  \item il codice diventa meno manutenibile e più difficile da capire.
\end{itemize}

\begin{esempio}{Bug che si scoprono solo a runtime}
\begin{affianco}[0.6][c]
\codicepiccolo
\begin{lstlisting}[style=js, name=bugs.js]
let message = "Hello Web Technologies!";
console.log(message.toLowerCase());   // hello web technologies!
message();       // TypeError: message is not a function (a runtime)

let employee = { name: "Jordan Belfort", role: "Stockbroker" };
employee.nome = "The Wolf of Wall Street";   // nessun errore!
console.log(employee.name);                  // Jordan Belfort

let team = ["Jordan Belfort", "Donnie Azoff", "Chester Ming"];
team.add("Nicky Koskoff");   // TypeError: team.add is not a function
\end{lstlisting}
\risultato
\begin{console}
hello web technologies!
Uncaught TypeError: message is
  not a function
\end{console}
\wtnota{L'esecuzione si ferma al primo errore;\\
        il refuso \code{nome} non ne produce affatto.}
\end{affianco}

Il secondo caso è il più insidioso: un refuso in un nome di proprietà
(\code{nome} invece di \code{name}) non produce \textbf{alcun} errore. Sarebbe
bello intercettare questi bug \textbf{prima} dell'esecuzione.
\end{esempio}

\section{I transpiler}

JavaScript ha cominciato a essere usato per programmi sempre più complessi, e
gli sviluppatori hanno sentito il bisogno di superare questi limiti. Sostituire
JavaScript con un linguaggio nuovo non era praticabile: era già supportato
ovunque.

Sono così nati molti linguaggi che \textbf{compilano verso JavaScript
standard}: CoffeeScript, TypeScript, Dart, Elm e circa altri 350.

\begin{figure}[H]
  \centering
  \begin{tikzpicture}[font=\Lato\small,
      n/.style={wtnode, minimum width=3.4cm, minimum height=1.3cm},
      a/.style={-{Stealth[length=6pt,width=5pt]}, line width=1pt,
                draw=neutral!45!black}]
    \node[n] (src) at (0,0)
      {codice in un nuovo\\ linguaggio, con\\ funzionalità in più};
    \node[n, draw=orange-gradient-6, fill=orange-gradient-1!40!white]
      (tr) at (5.0,0) {\textbf{Transpiler}};
    \node[n] (js) at (10.0,0)
      {codice JavaScript\\ (gira ovunque giri\\ JavaScript)};
    \draw[a] (src) -- (tr);
    \draw[a] (tr) -- (js);
  \end{tikzpicture}
  \caption{Un transpiler traduce da un linguaggio a un altro, non a codice
           macchina.}
  \label{fig:transpiler}
\end{figure}

\dfn{TypeScript}{
  \term{TypeScript} nasce nel 2012, progettato e sviluppato da Microsoft (lo
  sviluppatore principale è Anders Hejlsberg). È un \textbf{superset} di
  JavaScript con tipizzazione statica, interfacce, enumerazioni e generici, e
  \textbf{compila verso JavaScript}. È di gran lunga la variante di JavaScript
  più diffusa.}

Essendo un superset, \textbf{qualunque codice JavaScript valido è anche codice
TypeScript valido}: si può adottare gradualmente, un file alla volta.

I numeri confermano il primato. Nel sondaggio \textit{State of JavaScript 2022},
fra le 27\,294 persone che hanno risposto alla domanda sui linguaggi che
compilano verso JavaScript, 26\,988 hanno dichiarato di usare TypeScript; il
secondo classificato, Elm, si ferma a 624.

\begin{figure}[H]
  \centering
  \begin{tikzpicture}[font=\Lato\footnotesize]
    % lunghezza della barra = utenti / 3000 (in cm): pgfmath non accetta
    % numeri oltre 16383, quindi le lunghezze sono gia' calcolate
    \foreach \nome/\lung/\val [count=\i] in {TypeScript/8.996/26988,
        Elm/0.208/624, CoffeeScript/0.171/512, Flow/0.154/462,
        PureScript/0.107/322, ClojureScript/0.100/301, ReScript/0.083/248,
        Reason/0.070/211, Imba/0.035/106} {
      \pgfmathsetmacro{\yy}{-0.5*\i}
      \ifnum\i=1 \def\colore{primary!75!white}\else\def\colore{neutral!40!white}\fi
      \node[anchor=east] at (-0.15,\yy) {\nome};
      \fill[\colore] (0,\yy-0.16) rectangle (\lung,\yy+0.16);
      \node[anchor=west, font=\FiraCodeNL\scriptsize] at (\lung+0.06,\yy) {\val};
    }
    \draw[neutral!55!black, line width=0.6pt] (0,-0.2) -- (0,-4.8);
  \end{tikzpicture}
  \caption{Persone che usano i linguaggi che compilano verso JavaScript, secondo
           il sondaggio State of JavaScript 2022 (i dati delle slide).}
  \label{fig:ts-adozione}
\end{figure}

\section{Primi passi}

TypeScript si installa come pacchetto npm:

\begin{console}[Terminale]
$ npm install -g typescript

added 2 packages in 837ms
$ tsc --version
Version 7.0.2
\end{console}

Il comando rende disponibile globalmente il compilatore \code{tsc}. In
alternativa lo si può installare \textbf{nel progetto}, come dipendenza di
sviluppo, e lanciare con \code{npx}, che esegue i comandi dei pacchetti
presenti nella cartella \code{node\_modules} locale:

\begin{console}[Terminale]
$ npm install --save-dev typescript

added 2 packages in 623ms
$ npx tsc --version
Version 7.0.2
\end{console}

Così ogni progetto usa la versione del compilatore scritta nel proprio
\code{package.json}. I file TypeScript hanno tipicamente estensione \code{.ts}.

\info{Le versioni di questo capitolo}{
  Gli esempi sono stati eseguiti con TypeScript 7.0.2 (luglio 2026) e Node.js
  22.22.0, mentre le slide sono state preparate con TypeScript 5. Nel frattempo
  alcune impostazioni predefinite sono cambiate: dalla versione 6.0 (marzo
  2026) i \textbf{controlli severi} sono attivi fin dall'inizio e il JavaScript
  prodotto segue lo standard ES2025. Dove l'output è diverso da quello delle
  slide lo segnaliamo; per ritrovare il comportamento di allora basta
  l'opzione \code{-{}-strict false}. Con la versione 7, inoltre, il compilatore
  è diventato un eseguibile nativo, riscritto in Go: per questo npm installa
  \textbf{due} pacchetti, \code{typescript} e quello con l'eseguibile adatto al
  sistema (qui \code{@typescript/typescript-linux-x64}).}

\subsection{Il primo programma}

Il codice JavaScript standard è anche codice TypeScript valido: questa funzione
di saluto si può eseguire con Node come un normale file JavaScript.

\begin{lstlisting}[style=ts, name=hello.ts]
// Funzione di saluto di livello industriale, buona per ogni uso:
function greet(entity, date) {
  console.log(`Hello ${entity}, today is ${date}!`);
}

greet("Web Technologies");
\end{lstlisting}

\begin{console}[Terminale]
$ node hello.ts
Hello Web Technologies, today is undefined!
\end{console}

Abbiamo dimenticato il secondo argomento. Che cosa succede se compiliamo il
file con \code{tsc}?

\begin{console}[Terminale]
$ tsc hello.ts
hello.ts:2:16 - error TS7006: Parameter 'entity' implicitly has an 'any' type.

2 function greet(entity, date) {
                 ~~~~~~

hello.ts:2:24 - error TS7006: Parameter 'date' implicitly has an 'any' type.

2 function greet(entity, date) {
                         ~~~~

hello.ts:6:1 - error TS2554: Expected 2 arguments, but got 1.

6 greet("Web Technologies");
  ~~~~~

  hello.ts:2:24 - An argument for 'date' was not provided.
    2 function greet(entity, date) {
                             ~~~~


Found 3 errors in the same file, starting at: hello.ts:2

$ ls
hello.js
hello.ts
\end{console}

L'ultimo errore, \code{TS2554}, è proprio il nostro bug: la funzione si aspetta
due argomenti e ne riceve uno, e la nota sotto indica il parametro rimasto
senza valore. Il compilatore lo ha individuato \textbf{pur avendo scritto
soltanto JavaScript standard}, senza una sola annotazione di tipo.

Nelle slide compare soltanto questo errore. Gli altri due (\code{TS7006})
segnalano che i parametri \code{entity} e \code{date} hanno implicitamente il
tipo \code{any}, cioè «qualunque»: è uno dei controlli severi che TypeScript
6.0 ha reso predefiniti (ne parliamo nella sezione \textit{Livelli di
severità}, in fondo al capitolo). Spegnendoli si ottiene l'output delle slide:

\begin{console}[Terminale]
$ tsc --strict false hello.ts
hello.ts:6:1 - error TS2554: Expected 2 arguments, but got 1.

6 greet("Web Technologies");
  ~~~~~

  hello.ts:2:24 - An argument for 'date' was not provided.
    2 function greet(entity, date) {
                             ~~~~


Found 1 error in hello.ts:6

\end{console}

L'\code{ls} dopo la compilazione rivela anche un altro dettaglio: nonostante gli errori,
\code{tsc} ha generato \code{hello.js}. Per impostazione predefinita il
compilatore produce comunque il JavaScript: gli errori di tipo avvertono di ciò
che potrebbe andare storto, ma non bloccano la traduzione (a meno di attivare
l'opzione \code{noEmitOnError}).

\subsection{Correggere e compilare}

Correggiamo il bug passando anche la data:

\begin{lstlisting}[style=ts, name=hello.ts]
// Funzione di saluto di livello industriale, buona per ogni uso:
function greet(entity, date) {
  console.log(`Hello ${entity}, today is ${date}!`);
}

greet("Web Technologies", Date());
\end{lstlisting}

Come nelle slide, compiliamo senza i controlli severi (con i controlli attivi
\code{tsc} ripeterebbe i due \code{TS7006}, generando comunque il file):

\begin{console}[Terminale]
$ ls
hello.ts
$ tsc --strict false hello.ts
$ ls
hello.js
hello.ts
$ node hello.js
Hello Web Technologies, today is Tue Nov 17 2026 10:26:54 GMT+0000 (Coordinated Universal Time)!
\end{console}

Il compilatore ha \textbf{tradotto} \code{hello.ts} nel nuovo file
\code{hello.js}, che si esegue con Node come al solito. Ecco il codice
generato:

\begin{lstlisting}[style=js, name=hello.js]
"use strict";
// Funzione di saluto di livello industriale, buona per ogni uso:
function greet(entity, date) {
    console.log(`Hello ${entity}, today is ${date}!`);
}
greet("Web Technologies", Date());
\end{lstlisting}

Il compilatore conserva commenti e indentazione (portata a quattro spazi): il
codice tradotto sembra scritto da una persona. In testa compare la direttiva
\code{"use strict"}, che attiva la modalità rigorosa di JavaScript (capitolo
5): TypeScript 7 la aggiunge di serie, mentre TypeScript 5, come vedremo fra
un attimo, per impostazione predefinita non lo faceva.

\subsection{Downleveling}

Nelle slide il file generato è diverso: il template literal è diventato una
catena di chiamate a \code{concat}. Lo si ottiene ancora oggi facendo compilare
il file a TypeScript 5, che \code{npx} sa scaricare al volo (l'opzione
\code{-p} indica il pacchetto da usare, \code{-{}-yes} evita la domanda di
conferma):

\begin{console}[Terminale]
$ npx --yes -p typescript@5.9.3 tsc --version
Version 5.9.3
$ npx --yes -p typescript@5.9.3 tsc hello.ts
$ cat hello.js
// Funzione di saluto di livello industriale, buona per ogni uso:
function greet(entity, date) {
    console.log("Hello ".concat(entity, ", today is ").concat(date, "!"));
}
greet("Web Technologies", Date());
\end{console}

I template literal sono una funzionalità di ES6 (2015), e TypeScript 5 produceva
per impostazione predefinita codice ES5 (2009); fino alla versione 4.9 il
bersaglio predefinito era addirittura ES3, del 1999, quello citato nelle slide.
TypeScript sa \textbf{riscrivere} il codice da versioni più recenti del
linguaggio a versioni più vecchie: è il \term{downleveling}. Possiamo usare le
funzionalità moderne e far girare il risultato anche su browser antichi.

\begin{figure}[H]
  \centering
  \begin{tikzpicture}[font=\Lato\small,
      n/.style={wtnode, minimum width=2.9cm, minimum height=1.0cm},
      a/.style={-{Stealth[length=6pt,width=5pt]}, line width=1pt,
                draw=neutral!45!black}]
    \node[n] (es6) at (0,0) {codice ES6\\ \FiraCodeNL\scriptsize`Hello \$\{x\}`};
    \node[n, draw=orange-gradient-6, fill=orange-gradient-1!40!white]
      (tsc) at (4.2,0) {\textbf{tsc}};
    \node[n] (es5) at (8.4,0) {codice ES5\\ \FiraCodeNL\scriptsize"Hello ".concat(x)};
    \draw[a] (es6) -- (tsc);
    \draw[a] (tsc) -- (es5);
  \end{tikzpicture}
  \caption{Il downleveling: la stessa funzione, riscritta con i mezzi di una
           versione più vecchia del linguaggio.}
  \label{fig:downleveling}
\end{figure}

Il bersaglio si sceglie con l'opzione \code{-{}-target}: con
\code{tsc -{}-target es6 hello.ts} il template literal resta com'è. Con
TypeScript 7 i bersagli ES3 ed ES5 non esistono più: il minimo è ES6 (che si
può scrivere anche \code{es2015}) e, se non si indica nulla, il compilatore
produce ES2025. Per questo il nostro \code{hello.js} conservava il template
literal.

\begin{console}[Terminale]
$ tsc --target es3 hello.ts
error TS6046: Argument for '--target' option must be: 'es6', 'es2015', 'es2016', 'es2017', 'es2018', 'es2019', 'es2020', 'es2021', 'es2022', 'es2023', 'es2024', 'es2025', 'esnext'.
$ tsc --target es5 hello.ts
error TS5108: Option 'target=ES5' has been removed. Please remove it from your configuration.

Found 1 error.

\end{console}

Attenzione anche ai trattini. Nelle slide il comando è scritto
\code{tsc –target es6 hello.ts}, con un trattino tipografico (–) al posto dei
due trattini normali: copiato nel terminale non funziona, perché \code{tsc}
prende \code{–target} ed \code{es6} per due nomi di file da compilare.

% Il trattino tipografico (U+2013) e' fuori dalle tabelle di listings, che lo
% stamperebbe fuori posto ("path –'target'" invece di "path '–target'"): nel
% sorgente e' scritto "¬", che listings sostituisce con il trattino.
{\lstset{extendedchars=true, literate={¬}{{\textendash}}1}%
\begin{console}[Terminale]
$ tsc ¬target es6 hello.ts
error TS6231: Could not resolve the path 'es6' with the extensions: '.ts', '.tsx', '.d.ts', '.cts', '.d.cts', '.mts', '.d.mts'.
  The file is in the program because:
    Root file specified for compilation
error TS6231: Could not resolve the path '¬target' with the extensions: '.ts', '.tsx', '.d.ts', '.cts', '.d.cts', '.mts', '.d.mts'.
  The file is in the program because:
    Root file specified for compilation

Found 2 errors.

\end{console}
}

Il downleveling, naturalmente, resta: riguarda le funzionalità arrivate dopo
ES6. Due esempi sono di ES2020: l'optional chaining \code{a?.b} (capitolo 5) e
l'operatore \code{a ?? b}, che vale \code{b} quando \code{a} è \code{null} o
\code{undefined}. Il \code{?} dopo un nome di parametro o di proprietà, nella
prima riga, li rende facoltativi; lo riprenderemo fra poco.

\begin{lstlisting}[style=ts, name=down.ts]
function greet(user?: { name?: string }) {
  const name = user?.name ?? "anonymous";
  console.log(`Hello ${name}!`);
}

greet();
greet({ name: "Web Technologies" });
\end{lstlisting}

\begin{console}[Terminale]
$ tsc down.ts
$ cat down.js
"use strict";
function greet(user) {
    const name = user?.name ?? "anonymous";
    console.log(`Hello ${name}!`);
}
greet();
greet({ name: "Web Technologies" });
$ tsc --target es2019 down.ts
$ cat down.js
"use strict";
function greet(user) {
    var _a;
    const name = (_a = user === null || user === void 0 ? void 0 : user.name) !== null && _a !== void 0 ? _a : "anonymous";
    console.log(`Hello ${name}!`);
}
greet();
greet({ name: "Web Technologies" });
$ node down.js
Hello anonymous!
Hello Web Technologies!
\end{console}

Con il bersaglio predefinito il codice passa invariato (a parte le annotazioni
di tipo, che spariscono sempre). Chiedendo ES2019 i due operatori vengono
riscritti con confronti espliciti con \code{null} e \code{undefined}
(\code{void 0} è un modo sicuro di scrivere \code{undefined}), e la variabile
d'appoggio \code{\_a} evita di valutare due volte \code{user.name}. Il
comportamento non cambia.

\section{Dichiarare i tipi}

Finora abbiamo scritto soltanto JavaScript. TypeScript permette di dichiarare i
tipi di variabili, parametri e valori restituiti; i tipi primitivi più comuni
sono \code{string}, \code{number} e \code{boolean}.

\begin{lstlisting}[style=ts, name=primitive_types.ts]
let courseName: string = "Web Technologies";
let credits: number = 6;
let isGreat: boolean = true;

console.log(`
  Welcome to ${courseName}!
  Credits: ${credits} ETCS
  Status: ${isGreat ? "Great" : "Could be better"}
`);
\end{lstlisting}

A questo punto il file \textbf{non è più JavaScript valido}. Nelle slide,
eseguito con Node 22.14, produce un \code{SyntaxError} sui due punti della
prima annotazione. Con Node 22.22 lo stesso errore si ottiene solo spegnendo
una funzionalità di cui parliamo fra poco:

\begin{console}[Terminale]
$ node --no-experimental-strip-types primitive_types.ts
[cartella del progetto]/primitive_types.ts:1
let courseName: string = "Web Technologies";
              ^

SyntaxError: Unexpected token ':'
    at wrapSafe (node:internal/modules/cjs/loader:1638:18)
    ...

Node.js v22.22.0
\end{console}

Per eseguire l'esempio, il file va prima compilato:

\begin{console}[Terminale]
$ tsc primitive_types.ts
$ node primitive_types.js

  Welcome to Web Technologies!
  Credits: 6 ETCS
  Status: Great

\end{console}

\subsection{Node e la rimozione dei tipi}

Da Node 22.18 (luglio 2025) il comando \code{node} accetta direttamente i file
\code{.ts}: prima di eseguirli ne \textbf{cancella} le annotazioni di tipo
(\term{type stripping}) e manda in esecuzione il JavaScript che resta.

\begin{console}[Terminale]
$ node primitive_types.ts

  Welcome to Web Technologies!
  Credits: 6 ETCS
  Status: Great

\end{console}

\warning{Togliere i tipi non vuol dire controllarli}{
  Node non verifica nulla: rimuove le annotazioni e basta. Il programma della
  prossima sottosezione, che \code{tsc} rifiuta, con Node gira senza proteste.
  Il controllo dei tipi resta compito del compilatore (e dell'editor, che usa
  lo stesso motore). Inoltre la rimozione funziona solo per la sintassi che si
  può cancellare senza conseguenze: le enumerazioni, per esempio, non ne fanno
  parte (sezione \textit{Tipi letterali ed enumerazioni}).}

\subsection{Il controllo dei tipi}

Ora TypeScript può farci da guardia. Proviamo ad assegnare a \code{credits} una
stringa (rispetto alle slide abbiamo aggiunto il \code{console.log} finale, per
vedere che cosa fa Node):

\begin{lstlisting}[style=ts, name=primitive_types.ts]
let courseName: string = "Web Technologies";
let credits: number = 6;
let isGreat: boolean = true;

credits = "9 CFU";   // in JavaScript sarebbe lecito
console.log(credits);
\end{lstlisting}

\begin{console}[Terminale]
$ tsc primitive_types.ts
primitive_types.ts:5:1 - error TS2322: Type 'string' is not assignable to type 'number'.

5 credits = "9 CFU";   // in JavaScript sarebbe lecito
  ~~~~~~~


Found 1 error in primitive_types.ts:5

$ node primitive_types.ts
9 CFU
\end{console}

Il compilatore indica il punto dell'errore (riga 5, colonna 1), il suo codice
(\code{TS2322}) e un estratto del sorgente con la parte incriminata
sottolineata. Node invece, come annunciato, esegue il file e stampa
\code{9 CFU}.

\subsection{Array}

Il tipo di un array si indica con la sintassi \code{tipo[]}.

\begin{lstlisting}[style=ts, name=array.ts]
let langs: string[] = ["TypeScript", "JavaScript", "PHP", "Java"];
let primes: number[] = [2, 5, 7, 97, 829, 839];

langs.push({ name: "C#" });   // in JavaScript sarebbe lecito
\end{lstlisting}

\begin{console}[Terminale]
$ tsc array.ts
array.ts:4:12 - error TS2345: Argument of type '{ name: string; }' is not assignable to parameter of type 'string'.

4 langs.push({ name: "C#" });   // in JavaScript sarebbe lecito
             ~~~~~~~~~~~~~~


Found 1 error in array.ts:4

\end{console}

\subsection{Il tipo any}

TypeScript include un tipo speciale: \code{any}. Quando un valore è di tipo
\code{any} possiamo farci qualunque cosa e il compilatore non protesterà.

\begin{lstlisting}[style=ts, name=any.ts]
let x: any;

x = { name: "Angular", lang: "TypeScript" };
x.foo();              // compila, ma a runtime solleva TypeError

x = "Hello!";
let n: number = x;    // compila: una stringa finisce in un number
\end{lstlisting}

\begin{console}[Terminale]
$ tsc any.ts
$ node any.js
[cartella del progetto]/any.js:4
x.foo(); // compila, ma a runtime solleva TypeError
  ^

TypeError: x.foo is not a function
    at Object.<anonymous> ([cartella del progetto]/any.js:4:3)
    ...

Node.js v22.22.0
\end{console}

\warning{\code{any} disattiva i controlli}{
  Usare \code{any} equivale a rinunciare a tutto ciò per cui si è adottato
  TypeScript. È talvolta inevitabile (per esempio interfacciandosi con librerie
  non tipizzate), ma va considerato un'eccezione consapevole, non una scorciatoia
  di comodo.}

\subsection{Funzioni}

Si possono specificare i tipi sia degli ingressi sia del valore restituito.

\begin{lstlisting}[style=ts, name=functions.ts]
function greet(name: string): string {
  return `Hello ${name.toUpperCase()}`;
}

console.log(greet(42));
\end{lstlisting}

\begin{console}[Terminale]
$ tsc functions.ts
functions.ts:5:19 - error TS2345: Argument of type 'number' is not assignable to parameter of type 'string'.

5 console.log(greet(42));
                    ~~


Found 1 error in functions.ts:5

\end{console}

\subsection{Inferenza automatica}

Quando non si fornisce un'annotazione esplicita, TypeScript prova a
\textbf{dedurre} il tipo dal contesto.

\begin{affianco}[0.62][c]
\codicepiccolo
\begin{lstlisting}[style=ts, name=inference.ts]
function greet(name) {             // function greet(name: any): string
  return `Hello ${name}`;
}

let arr = [1, 2, 3, 4, 5];         // let arr: number[]

let course = "Web Technologies";   // let course: string

greet(course);

course = { name: "Web Technologies", credits: 6 };
\end{lstlisting}
\risultato
\begin{immagine}[Ciò che mostra l'editor passando sopra\\ ai nomi: i tipi dedotti.]
\begin{tikzpicture}[
    tip/.style={fill=consolebg, text=consolefg, rounded corners=2pt,
      font=\FiraCodeNL\scriptsize, inner sep=3pt, anchor=west}]
  \node[tip] at (0,0)    {function greet(name: any): string};
  \node[tip] at (0,-0.6) {let arr: number[]};
  \node[tip] at (0,-1.2) {let course: string};
\end{tikzpicture}
\end{immagine}
\end{affianco}

Il tipo dedotto per \code{course} è \code{string}, e l'ultimo assegnamento lo
viola. Con TypeScript 7 gli errori sono però due: anche il parametro
\code{name} viene segnalato, perché il suo tipo non si può dedurre e resta un
\code{any} implicito (il tipo restituito invece sì: \code{string}). Senza
controlli severi resta solo l'errore delle slide:

\begin{console}[Terminale]
$ tsc inference.ts
inference.ts:1:16 - error TS7006: Parameter 'name' implicitly has an 'any' type.

1 function greet(name) {             // function greet(name: any): string
                 ~~~~

inference.ts:11:1 - error TS2322: Type '{ name: string; credits: number; }' is not assignable to type 'string'.

11 course = { name: "Web Technologies", credits: 6 };
   ~~~~~~


Found 2 errors in the same file, starting at: inference.ts:1

$ tsc --strict false inference.ts
inference.ts:11:1 - error TS2322: Type '{ name: string; credits: number; }' is not assignable to type 'string'.

11 course = { name: "Web Technologies", credits: 6 };
   ~~~~~~


Found 1 error in inference.ts:11

\end{console}

\subsection{Gli any impliciti}

Quando non si specifica un tipo e TypeScript non riesce a dedurlo, il
compilatore ripiega su \code{any}. Conviene evitarlo, perché su un \code{any}
non si fa alcun controllo: il flag \code{noImplicitAny} segnala come errore
ogni \code{any} implicito.

\begin{lstlisting}[style=ts, name=implicit_any.ts]
let x;                    // let x: any

function greet(name) {    // function greet(name: any): void
  console.log(`Hello, ${name}`);
}

greet(x);
\end{lstlisting}

\begin{console}[Terminale]
$ tsc --strict false implicit_any.ts
$ tsc --strict false --noImplicitAny implicit_any.ts
implicit_any.ts:3:16 - error TS7006: Parameter 'name' implicitly has an 'any' type.

3 function greet(name) {    // function greet(name: any): void
                 ~~~~


Found 1 error in implicit_any.ts:3

$ tsc implicit_any.ts
implicit_any.ts:3:16 - error TS7006: Parameter 'name' implicitly has an 'any' type.

3 function greet(name) {    // function greet(name: any): void
                 ~~~~


Found 1 error in implicit_any.ts:3

\end{console}

I primi due comandi riproducono le slide, dove i controlli severi sono spenti:
senza flag il file compila, con \code{-{}-noImplicitAny} no. Con TypeScript 6 e
7 il flag è già attivo, perché fa parte della modalità severa, come mostra
l'ultimo comando. La riga \code{let x;} invece non viene segnalata: a una
variabile dichiarata senza valore TypeScript attribuisce, punto per punto, il
tipo dell'ultimo valore ricevuto (al momento della chiamata, \code{undefined}).

\section{Tipi oggetto}

Un tipo oggetto si definisce elencando le sue proprietà e i loro tipi.

\begin{affianco}[0.6][c]
\codicepiccolo
\begin{lstlisting}[style=ts, name=object_types.ts]
function printCourse(course: { name: string, credits: number }) {
  console.log(`Course name: ${course.name}`);
  console.log(`Course credits: ${course.credits}`);
}

printCourse({ name: "Web Technologies", credits: 6 });   // ok

printCourse({ name: "Software Engineering" });   // errore

printCourse({ name: "Algebra", credits: 6, year: "First" });   // errore
\end{lstlisting}
\risultato
\begin{immagine}[L'oggetto deve avere esattamente la forma attesa.]
\begin{tikzpicture}[
    f/.style={rounded corners=2pt, draw=neutral!55!black, fill=white, line width=0.6pt,
      font=\FiraCodeNL\scriptsize, inner sep=3pt, anchor=west}]
  \node[f, draw=primary!60!black, fill=primary!10!white] (t) at (0,0) {\{ name: string, credits: number \}};
  \node[mnota, anchor=south west] at (t.north west) {atteso};
  \node[f] (a) at (0,-0.75) {\{ name, credits \}};   \node[anchor=west] at ([xshift=3pt]a.east) {\mok};
  \node[f] (b) at (0,-1.35) {\{ name \}};            \node[anchor=west] (bx) at ([xshift=3pt]b.east) {\merr};
  \node[mnota, anchor=west] at (bx.east) {manca \code{credits}};
  \node[f] (c) at (0,-1.95) {\{ name, credits, year \}}; \node[anchor=west] (cx) at ([xshift=3pt]c.east) {\merr};
  \node[mnota, anchor=west] at (cx.east) {\code{year} in più};
\end{tikzpicture}
\end{immagine}
\end{affianco}

\begin{console}[Terminale]
$ tsc object_types.ts
object_types.ts:8:13 - error TS2741: Property 'credits' is missing in type '{ name: string; }' but required in type '{ name: string; credits: number; }'.

8 printCourse({ name: "Software Engineering" });   // errore
              ~~~~~~~~~~~~~~~~~~~~~~~~~~~~~~~~

  object_types.ts:1:46 - 'credits' is declared here.
    1 function printCourse(course: { name: string, credits: number }) {
                                                   ~~~~~~~

object_types.ts:10:44 - error TS2353: Object literal may only specify known properties, and 'year' does not exist in type '{ name: string; credits: number; }'.

10 printCourse({ name: "Algebra", credits: 6, year: "First" });   // errore
                                              ~~~~


Found 2 errors in the same file, starting at: object_types.ts:8

\end{console}

Nelle slide il primo errore porta il codice TS2453, ma è un'inversione di
cifre: TypeScript 5 lo segnala come \code{TS2345} (\textit{Argument of type
\ldots{} is not assignable to parameter of type \ldots}), seguito dalla riga
sulla proprietà mancante. TypeScript 7 riporta direttamente quest'ultima, con il
codice \code{TS2741}, e indica dove \code{credits} è dichiarata.

\subsection{Proprietà opzionali}

Una proprietà si può dichiarare \textbf{opzionale} aggiungendo un \code{?} dopo
il nome.

\begin{lstlisting}[style=ts, name=optional.ts]
function printCourse(course: { name: string, credits?: number }) {
  console.log(`Course name: ${course.name}`);
  console.log(`Course credits: ${course.credits}`);
  // andrebbe verificato se course.credits e' undefined!
}

printCourse({ name: "Web Technologies", credits: 6 });   // ok
printCourse({ name: "Software Engineering" });   // ok, ma...

printCourse({ name: "Algebra", credits: 6, year: "First" });   // errore
\end{lstlisting}

\begin{console}[Terminale]
$ tsc optional.ts
optional.ts:10:44 - error TS2353: Object literal may only specify known properties, and 'year' does not exist in type '{ name: string; credits?: number | undefined; }'.

10 printCourse({ name: "Algebra", credits: 6, year: "First" });   // errore
                                              ~~~~


Found 1 error in optional.ts:10

\end{console}

Il messaggio riporta il tipo con la proprietà opzionale, \code{credits?: number
| undefined} (nelle slide il commento ripete per svista il tipo dell'esempio
precedente; il \code{| undefined} compare con i controlli severi). Commentata
l'ultima chiamata, il programma compila e mostra il problema annunciato nel
commento:

\begin{console}[Terminale]
$ tsc optional.ts
$ node optional.js
Course name: Web Technologies
Course credits: 6
Course name: Software Engineering
Course credits: undefined
\end{console}

\section{Tipi unione}

I \term{tipi unione} si ottengono combinando due o più tipi, elencati e separati
da \code{|}. Rappresentano valori che possono appartenere a uno qualunque dei
tipi combinati.

\begin{lstlisting}[style=ts, name=union.ts]
function printError(code: number | string) {
  console.log(`Error code: ${code}`);
}

printError(404);                 // ok
printError("401 Unauthorized");  // ok
printError({ code: "403", message: "Forbidden" });   // errore
\end{lstlisting}

\begin{console}[Terminale]
$ tsc union.ts
union.ts:7:12 - error TS2345: Argument of type '{ code: string; message: string; }' is not assignable to parameter of type 'string | number'.

7 printError({ code: "403", message: "Forbidden" });   // errore
             ~~~~~~~~~~~~~~~~~~~~~~~~~~~~~~~~~~~~~


Found 1 error in union.ts:7

\end{console}

TypeScript garantisce che le operazioni su un tipo unione siano valide per
\textbf{ciascuno} dei tipi possibili:

\begin{lstlisting}[style=ts, name=union2.ts]
function printError(code: number | string) {
  console.log(`Error code: ${code.toUpperCase()}`);   // errore
}
\end{lstlisting}

\begin{console}[Terminale]
$ tsc union2.ts
union2.ts:2:35 - error TS2339: Property 'toUpperCase' does not exist on type 'string | number'.
  Property 'toUpperCase' does not exist on type 'number'.

2   console.log(`Error code: ${code.toUpperCase()}`);   // errore
                                    ~~~~~~~~~~~


Found 1 error in union2.ts:2

\end{console}

In questi casi bisogna \term{restringere} il tipo con controlli espliciti (per
esempio \code{typeof} o \code{Array.isArray()}); TypeScript sa dedurre tipi
specifici nei singoli rami.

\begin{affianco}[0.6][c]
\codicepiccolo
\begin{lstlisting}[style=ts, name=narrowing.ts]
function printError(code: number | string) {
  if (typeof code === "string") {
    console.log(`Error: ${code.toUpperCase()}`);   // qui code: string
  } else {
    console.log(`Error: ${code}`);                 // qui code: number
  }
}

printError("not found");
printError(404);
\end{lstlisting}
\risultato
\begin{immagine}[Dentro ogni ramo il tipo è più preciso.]
\begin{tikzpicture}[
    t/.style={rounded corners=2pt, draw=primary!60!black, fill=primary!10!white,
      line width=0.7pt, font=\FiraCodeNL\scriptsize, inner sep=3pt},
    a/.style={-{Stealth[length=4pt,width=3.5pt]}, line width=0.7pt, draw=neutral!45!black}]
  \node[t] (u) at (0,0) {code: number | string};
  \node[t, fill=green-gradient-1!45!white, draw=green-gradient-6] (s) at (-1.3,-1.4) {code: string};
  \node[t, fill=accent!12!white, draw=accent!55!black] (n) at (1.4,-1.4) {code: number};
  \draw[a] (u) -- node[mnota, left, align=right]{\code{typeof}\\ \code{=== "string"}} (s);
  \draw[a] (u) -- node[mnota, right]{altrimenti} (n);
  \node[mnota, anchor=north] at (s.south) {\code{toUpperCase()} \mok};
\end{tikzpicture}
\end{immagine}
\end{affianco}

\begin{console}[Terminale]
$ tsc narrowing.ts
$ node narrowing.js
Error: NOT FOUND
Error: 404
\end{console}

\section{Alias di tipo e interfacce}

Finora abbiamo scritto i tipi oggetto e i tipi unione direttamente nelle
annotazioni. Volendo riusare lo stesso tipo più volte (la priorità è scrivere
codice \term{DRY}, \textit{Don't Repeat Yourself}), gli si può assegnare un
nome: è un \term{alias di tipo}, introdotto dalla parola chiave \code{type}
con la sintassi \code{type <nome> = <definizione>}.

\begin{affianco}[0.6][c]
\codicepiccolo
\begin{lstlisting}[style=ts, name=aliases.ts]
type Credits = number | string;

type Course = {
  name: string,
  credits: Credits
}

function printCourse(course: Course) {
  console.log(`Name: ${course.name}, Credits: ${course.credits}`);
}

let webtech: Course = { name: "Web Technologies", credits: 9 };
let softeng: Course = { name: "Software Engineering", credits: "9 CFU" };

printCourse(webtech);
printCourse(softeng);
\end{lstlisting}
\risultato
\begin{immagine}[Un alias può comporre altri alias.]
\begin{tikzpicture}[    f/.style={rounded corners=2pt, draw=primary!60!black, fill=primary!8!white,
      line width=0.7pt, font=\FiraCodeNL\scriptsize, inner sep=3pt, align=left},
    a/.style={-{Stealth[length=4pt,width=3.5pt]}, line width=0.7pt, draw=neutral!45!black}]
  \node[f] (c) at (0,0) {type Course\\ \ \ name: string\\ \ \ credits: \textcolor{accent!70!black}{Credits}};
  \node[f, fill=accent!12!white, draw=accent!55!black] (k) at (3.3,-0.35) {type Credits\\ \ \ number | string};
  \draw[a] ([yshift=-0.35cm]c.east) -- (k.west);
\end{tikzpicture}
\end{immagine}
\end{affianco}

\begin{console}[Terminale]
$ tsc aliases.ts
$ node aliases.js
Name: Web Technologies, Credits: 9
Name: Software Engineering, Credits: 9 CFU
\end{console}

Le \term{interfacce} sono un altro modo di dare un nome ai tipi oggetto.

\begin{lstlisting}[style=ts, name=interfaces.ts]
interface Course {
  name: string;
  credits: number;
}

function printCourse(course: Course) {
  console.log(`Name: ${course.name}, Credits: ${course.credits}`);
}

let webtech: Course = { name: "Web Technologies", credits: 9 };

printCourse(webtech);
printCourse({ name: "Software Engineering", credits: 9 });
\end{lstlisting}

\begin{console}[Terminale]
$ tsc interfaces.ts
$ node interfaces.js
Name: Web Technologies, Credits: 9
Name: Software Engineering, Credits: 9
\end{console}

\subsection{Tipizzazione strutturale}

\dfn{Tipizzazione strutturale}{
  TypeScript è un linguaggio a \term{tipizzazione strutturale}: per stabilire la
  compatibilità fra tipi guarda soltanto alla \textbf{struttura} e alle
  capacità, non ai \textbf{nomi}. Java, al contrario, ha un sistema di tipi
  \textit{nominativo}.}

\begin{affianco}[0.55][c]
\begin{lstlisting}[style=ts, name=structural.ts]
interface Person { name: string; age: number; }
interface Pet    { name: string; age: number; }

function printPerson(p: Person) {
  console.log(`Name: ${p.name}, age: ${p.age}`);
}

let person: Person = { name: "Janet", age: 20 };
let pet: Pet       = { name: "Chuck", age: 3 };

printPerson(person);   // Name: Janet, age: 20
printPerson(pet);      // Name: Chuck, age: 3   (!)
\end{lstlisting}
\risultato
\begin{immagine}[Stessa forma, quindi stesso tipo:\\ i nomi non contano.]
\begin{tikzpicture}[    f/.style={rounded corners=2pt, draw=primary!60!black, fill=primary!8!white,
      line width=0.7pt, font=\FiraCodeNL\scriptsize, inner sep=3pt, align=left},
    a/.style={-{Stealth[length=4pt,width=3.5pt]}, line width=0.7pt, draw=neutral!45!black}]
  \node[f] (p) at (0,0) {Person\\ \ \ name: string\\ \ \ age: number};
  \node[f, fill=accent!12!white, draw=accent!55!black] (q) at (3.0,0) {Pet\\ \ \ name: string\\ \ \ age: number};
  \node at (1.5,0) {\Large =};
  \node[mnota, anchor=north] at (1.5,-0.8) {TypeScript \mok\quad Java \merr};
\end{tikzpicture}
\end{immagine}
\end{affianco}

Un \code{Pet} viene accettato dove è atteso un \code{Person} perché ne ha la
stessa forma: il file compila e stampa le due righe dei commenti. In Java
sarebbe un errore di compilazione.

\subsection{Estendere interfacce e tipi}

Un'interfaccia può \textbf{estenderne} un'altra con \code{extends}: eredita
tutte le sue proprietà e ne aggiunge di nuove.

\begin{affianco}[0.5][c]
\begin{lstlisting}[style=ts, name=extending.ts]
interface Pet {
  name: string
}

interface Bird extends Pet {
  flies: boolean
}

let chuck: Bird = { name: "Chuck", flies: true };
let quentin: Bird = { flies: false };   // errore
\end{lstlisting}
\risultato
\begin{immagine}[\code{Bird} ha tutto ciò che ha \code{Pet}, più \code{flies}.]
\begin{tikzpicture}[
    r/.style={draw=primary!60!black, fill=primary!8!white, line width=0.7pt,
      font=\FiraCodeNL\scriptsize, minimum width=2.6cm, minimum height=0.45cm},
    n/.style={r, fill=accent!12!white, draw=accent!55!black}]
  \node[r] (n1) at (0,0) {name: string};
  \node[n, anchor=north] (n2) at ([yshift=0.7pt]n1.south) {flies: boolean};
  \draw[wtbrace] ([xshift=-3pt]n1.north west) -- ([xshift=-3pt]n1.south west)
    node[midway, left=4pt, mnota]{Pet};
  \draw[wtbrace, decoration={mirror=false}] ([xshift=3pt]n1.north east) -- ([xshift=3pt]n2.south east)
    node[midway, right=4pt, mnota]{Bird};
\end{tikzpicture}
\end{immagine}
\end{affianco}

\begin{console}[Terminale]
$ tsc extending.ts
extending.ts:10:5 - error TS2741: Property 'name' is missing in type '{ flies: false; }' but required in type 'Bird'.

10 let quentin: Bird = { flies: false };   // errore
       ~~~~~~~

  extending.ts:2:3 - 'name' is declared here.
    2   name: string
        ~~~~


Found 1 error in extending.ts:10

\end{console}

Nelle slide l'errore è un \code{TS2322} che dice \textit{required in type
'Pet'}. Le versioni recenti (già la 5.9) segnalano invece \code{TS2741}, citano
\code{Bird}, il tipo della variabile, e indicano in \code{Pet} la dichiarazione
della proprietà mancante. La sostanza non cambia: a \code{quentin} manca il
nome.

Alias di tipo e interfacce sono molto simili e spesso intercambiabili. Una
distinzione chiave: i \code{type} \textbf{non possono essere riaperti} per
aggiungere proprietà, mentre le interfacce sono sempre estendibili. Ecco lo
stesso codice nelle due versioni:

\begin{affianco}[0.5]
\codicepiccolo
\begin{lstlisting}[style=ts, name=reopen_interface.ts]
interface Pet {
  name: string
}

interface Pet {
  age: number
}

let matt: Pet = { name: "Matt", age: 2 };
\end{lstlisting}
\risultato
\codicepiccolo
\begin{lstlisting}[style=ts, name=reopen_type.ts]
type Pet = {
  name: string
}

type Pet = {
  age: number
}

let matt: Pet = { name: "Matt", age: 2 };
\end{lstlisting}
\end{affianco}

\begin{console}[Terminale]
$ tsc reopen_interface.ts
$ tsc reopen_type.ts
reopen_type.ts:1:6 - error TS2300: Duplicate identifier 'Pet'.

1 type Pet = {
       ~~~

reopen_type.ts:5:6 - error TS2300: Duplicate identifier 'Pet'.

5 type Pet = {
       ~~~

reopen_type.ts:9:33 - error TS2353: Object literal may only specify known properties, and 'age' does not exist in type 'Pet'.

9 let matt: Pet = { name: "Matt", age: 2 };
                                  ~~~


Found 3 errors in the same file, starting at: reopen_type.ts:1

\end{console}

Le due dichiarazioni dell'interfaccia si fondono in una sola, con entrambe le
proprietà. Con \code{type}, invece, il nome risulta dichiarato due volte
(\code{TS2300}, segnalato su tutte e due le dichiarazioni) e il compilatore
tiene la prima: per questo \code{age} diventa una proprietà sconosciuta.

Gli alias si estendono con i \term{tipi intersezione}, creati con \code{\&}.
Come le unioni, le intersezioni combinano più tipi in uno; il tipo risultante
ha le proprietà di \textbf{ciascuno} dei tipi combinati.

\begin{lstlisting}[style=ts, name=intersection.ts]
type Pet = { name: string }

type Bird = Pet & {
  flies: boolean
}

let chuck: Bird = { name: "Chuck", flies: true };
\end{lstlisting}

\section{Asserzioni di tipo}

Talvolta lo sviluppatore ne sa più del compilatore. Per esempio, usando
\code{document.getElementById()} TypeScript può dedurre soltanto che la chiamata
restituirà un \code{HTMLElement}, mentre noi sappiamo che quell'elemento è di un
tipo più specifico.

\begin{lstlisting}[style=ts]
const nameInput = document.getElementById("name") as HTMLInputElement;
\end{lstlisting}

\warning{Non è un'asserzione nel senso dei test}{
  Il termine è fuorviante: non viene verificato nulla a runtime. È soltanto
  un'indicazione al compilatore, valida a tempo di compilazione. Se ci
  sbagliamo, l'errore si manifesterà comunque durante l'esecuzione.}

\section{Tipi letterali ed enumerazioni}

Oltre ai tipi \code{string} e \code{number}, si può fare riferimento a stringhe
e numeri \textbf{specifici} in posizione di tipo.

\begin{lstlisting}[style=ts, name=literals.ts]
let alignment: "center";
alignment = "center";
alignment = "left";
// TS2322: Type '"left"' is not assignable to type '"center"'
\end{lstlisting}

Da solo non è molto utile, ma combinato con i tipi unione esprime un concetto
molto potente: tipi che corrispondono a un \textbf{insieme di valori noti}.

\begin{lstlisting}[style=ts, name=literals2.ts]
type Alignment = "center" | "left" | "right";

let alignment: Alignment;
alignment = "left";    // ok
alignment = "right";   // ok
alignment = "centre";
// TS2820: Type '"centre"' is not assignable to type 'Alignment'.
// Did you mean '"center"'?
\end{lstlisting}

Le \term{enumerazioni} permettono di definire un insieme di costanti con un
nome.

\begin{affianco}[0.55][c]
\begin{lstlisting}[style=ts, name=enums.ts]
interface Order { isPremium: boolean }

enum Priority {
  Low,
  Medium,
  High
}

function computePriority(s: Order): Priority {
  if (s.isPremium)
    return Priority.High;
  return Priority.Low;
}

console.log(Priority.Medium);   // 1
\end{lstlisting}
\risultato
\begin{immagine}[Senza valori espliciti, le costanti\\ sono numerate da 0.]
\begin{tikzpicture}
  \mchiave{1.4cm}
  \oggetto{e}{(0,0)}{2.4cm}{enum Priority}{Low/0, Medium/1, High/2}
\end{tikzpicture}
\end{immagine}
\end{affianco}

A differenza dei costrutti visti finora, un'enumerazione \textbf{non sparisce}
nella traduzione: diventa un oggetto JavaScript, con la corrispondenza nei due
sensi (\code{Priority.Medium} vale \code{1} e \code{Priority[1]} vale
\code{"Medium"}).

\begin{console}[Terminale]
$ tsc enums.ts
$ node enums.js
1
$ cat enums.js
"use strict";
var Priority;
(function (Priority) {
    Priority[Priority["Low"] = 0] = "Low";
    Priority[Priority["Medium"] = 1] = "Medium";
    Priority[Priority["High"] = 2] = "High";
})(Priority || (Priority = {}));
function computePriority(s) {
    if (s.isPremium)
        return Priority.High;
    return Priority.Low;
}
console.log(Priority.Medium); // 1
\end{console}

Per questo Node non può eseguire il file limitandosi a togliere i tipi:

\begin{console}[Terminale]
$ node enums.ts
[cartella del progetto]/enums.ts:3

  > enum Priority {
      Low,
      Medium,
      High
  > }

SyntaxError [ERR_UNSUPPORTED_TYPESCRIPT_SYNTAX]: TypeScript enum is not supported in strip-only mode
    at parseTypeScript (node:internal/modules/typescript:63:40)
    ...

Node.js v22.22.0
\end{console}

\section{Tipi funzione}

Come si dichiara che una funzione riceve una \textbf{callback} con certi
parametri e un certo risultato? Con le \term{espressioni di tipo funzione}.

\begin{lstlisting}[style=ts, name=function_types.ts]
type stringDecoratorFunction = (x: string, mode: boolean) => string;

let fn: () => string;                  // funzione senza argomenti, ritorna string
fn = () => { return "Hello Web Tech" };       // ok
fn = (x: string) => { return `Hello ${x}` };  // errore di tipo
\end{lstlisting}

Il tipo \code{stringDecoratorFunction} descrive una funzione che riceve una
\code{string} e un \code{boolean} e restituisce una \code{string}. La variabile
\code{fn} invece può contenere soltanto funzioni senza argomenti: la seconda
funzione ne vuole uno, che chi chiama \code{fn} non passerà mai.

\begin{console}[Terminale]
$ tsc function_types.ts
function_types.ts:5:1 - error TS2322: Type '(x: string) => string' is not assignable to type '() => string'.
  Target signature provides too few arguments. Expected 1 or more, but got 0.

5 fn = (x: string) => { return `Hello ${x}` };  // errore di tipo
  ~~


Found 1 error in function_types.ts:5

\end{console}

Nelle slide l'ultima riga è commentata come \textit{TypeError}: non si tratta
dell'eccezione di JavaScript a runtime, ma di un errore di tipo segnalato dal
compilatore (\code{TS2322}), prima di eseguire alcunché.

Un tipo funzione si può scrivere anche direttamente nella firma, come per il
parametro \code{decorator} qui sotto:

\begin{affianco}[0.62][c]
\codicepiccolo
\begin{lstlisting}[style=ts, name=greet_case.ts]
enum Case { Uppercase, Lowercase }

function GREET(name: string, decorator: (x: string, mode: Case) => string) {
  return decorator(name, Case.Uppercase);
}

function stringDecorator(x: string, mode: Case) {
  if (mode === Case.Lowercase)
    return `>>> Hello ${x} <<<`.toLowerCase();
  else
    return `>>> Hello ${x} <<<`.toUpperCase();
}

console.log(GREET("Web Technologies", stringDecorator));
\end{lstlisting}
\risultato
\begin{console}[Terminale]
$ tsc greet_case.ts
$ node greet_case.js
>>> HELLO WEB TECHNOLOGIES <<<
\end{console}
\wtnota{\code{GREET} chiama la callback\\ con \code{Case.Uppercase}.}
\end{affianco}

\section{Generici}

Buona parte dell'ingegneria del software mira a costruire software
\textbf{riusabile}: sviluppare componenti capaci di operare su dati presenti e
futuri è essenziale per i sistemi di grandi dimensioni. I \term{generici} sono
uno degli strumenti principali per scrivere codice che funzioni con una varietà
di tipi.

\begin{esempio}{La funzione identità}
Supponiamo di dover realizzare una funzione che restituisce ciò che riceve.

\begin{lstlisting}[style=ts, name=identity_any.ts]
function identity(x) {        // function identity(x: any): any
  return x;
}

let y: string = "Hello";
let z = identity(y);          // z: any (!)
\end{lstlisting}

Usare \code{any} (implicitamente o esplicitamente) è abbastanza generico, ma
\textbf{perdiamo l'informazione} sul tipo al ritorno. (Con TypeScript 7 questa
versione non compila nemmeno: \code{x} è un \code{any} implicito, errore
\code{TS7006}.) Dare un tipo specifico risolve il problema, ma la funzione
servirebbe solo per le stringhe:

\begin{lstlisting}[style=ts, name=identity_string.ts]
function identity(x: string): string {
  return x;
}

let y: string = "Hello";
let z = identity(y);          // z: string (ma funziona solo con le stringhe)
\end{lstlisting}

TypeScript permette di definire \term{variabili di tipo}, che operano sui tipi
invece che sui valori. Si dichiarano dopo il nome della funzione, fra \code{<}
e \code{>}; qui la variabile si chiama \code{Type}, ma va bene qualunque
identificatore.

\begin{affianco}[0.55][c]
\codicepiccolo
\begin{lstlisting}[style=ts, name=generics.ts]
function identity<Type>(x: Type): Type {
  return x;
}

let s: string = identity<string>("Hello");   // Type esplicito
let n: number = identity(42);                // dedotto automaticamente
let o: { title: string, artist: string } = identity({
  title: "Sultans of swing",
  artist: "Dire Straits"
});
\end{lstlisting}
\risultato
\begin{immagine}[Il tipo in uscita è quello in ingresso.]
\begin{tikzpicture}[
    t/.style={rounded corners=2pt, draw=primary!60!black, fill=primary!10!white,
      line width=0.7pt, font=\FiraCodeNL\scriptsize, inner sep=3pt},
    fn/.style={mfun},
    a/.style={-{Stealth[length=4pt,width=3.5pt]}, line width=0.7pt, draw=neutral!45!black}]
  \foreach \t/\y in {string/0, number/-0.7, {\{ title, artist \}}/-1.4} {
    \node[t, anchor=east] (i) at (-0.1,\y) {\t};
    \node[t, anchor=west] (o) at (2.5,\y) {\t};
    \draw[a] (i.east) -- (0.35,\y); \draw[a] (2.05,\y) -- (o.west);
  }
  \node[fn, minimum height=2.0cm, minimum width=1.6cm] at (1.2,-0.7) {identity};
\end{tikzpicture}
\end{immagine}
\end{affianco}

Non è la stessa cosa di usare \code{any}: qui l'informazione sul tipo di
ingresso viene \textbf{preservata}.
\end{esempio}

Le variabili di tipo possono essere più di una, e funzionano anche con gli
array:

\begin{lstlisting}[style=ts, name=merge.ts]
function merge<T1, T2>(x: T1, y: T2): { field: T2, otherField: T1 } {
  return { field: y, otherField: x };
}

let x = merge("Sultans of Swing", 42);
// x ha tipo { field: number, otherField: string }

function getFirst<T>(arr: T[]): T {
  return arr[0];
}

let first: string = getFirst(["hello", "web", "technologies"]);
console.log(first);   // hello
\end{lstlisting}

\subsection{Classi generiche}

Anche una classe può avere variabili di tipo. Questa collezione restituisce un
elemento a caso, del tipo scelto alla creazione:

\begin{lstlisting}[style=ts, name=generic_class.ts]
class CustomCollection<T> {
  list: T[] = [];

  addElement(element: T): number {
    return this.list.push(element);
  }

  getRandomElement(): T {
    let selectedIndex: number = Math.floor(Math.random() * this.list.length);
    return this.list[selectedIndex];
  }
}

let c = new CustomCollection<string>();
c.addElement("hello"); c.addElement("web"); c.addElement("technologies");
console.log(c.getRandomElement());
\end{lstlisting}

\begin{console}[Terminale]
$ tsc generic_class.ts
$ for i in 1 2 3 4 5; do node generic_class.js; done
web
technologies
technologies
hello
web
\end{console}

\subsection{Vincoli sui generici}

Con \code{extends} si può \textbf{vincolare} una variabile di tipo: qui
\code{T} può essere soltanto un \code{Animal} o un suo sottotipo, e per questo
dentro la funzione si può leggere \code{species}.

\begin{affianco}[0.6][c]
\codicepiccolo
\begin{lstlisting}[style=ts, name=constraints.ts]
class Animal { species: string; }
class Bird extends Animal  { canFly: boolean; }
class Snake extends Animal { isVenomous: boolean; }

function animalInfo<T extends Animal>(animal: T): void {
  console.log(animal.species);
}

let a: Animal = new Animal(); a.species = "Dog";
let b: Bird = new Bird();     b.species = "Kiwi";  b.canFly = false;
let c: Snake = new Snake();   c.species = "Cobra"; c.isVenomous = true;

animalInfo(a);                       // Dog
animalInfo(b);                       // Kiwi
animalInfo(c);                       // Cobra
animalInfo({ species: "Spider" });   // Spider (!)
\end{lstlisting}
\risultato
\begin{immagine}[Basta avere \code{species}: per la tipizzazione\\
                 strutturale anche l'oggetto letterale è un \code{Animal}.]
\begin{tikzpicture}[
    c/.style={rounded corners=2pt, draw=primary!60!black, fill=primary!8!white,
      line width=0.7pt, font=\FiraCodeNL\scriptsize, inner sep=3pt},
    a/.style={-{Stealth[length=5pt,width=4pt,open]}, line width=0.7pt, draw=neutral!45!black}]
  \node[c] (an) at (1.4,0) {Animal\\[-1pt] \scriptsize species};
  \node[c] (b) at (0,-1.4) {Bird};
  \node[c] (s) at (1.4,-1.4) {Snake};
  \node[c, dashed, fill=white] (l) at (3.1,-1.4) {\{ species \}};
  \draw[a] (b) -- (an); \draw[a] (s) -- (an);
  \draw[a, dashed] (l) -- (an);
\end{tikzpicture}
\end{immagine}
\end{affianco}

L'ultima riga funziona per via della tipizzazione strutturale: l'oggetto
letterale ha la stessa forma di \code{Animal}, quindi soddisfa il vincolo.

Con TypeScript 7, però, il file così com'è non compila. Le classi dichiarano
proprietà senza valore iniziale, e la modalità severa lo segnala
(\code{TS2564}): leggendo \code{species} prima di averla assegnata si
otterrebbe \code{undefined}, non una stringa.

\begin{console}[Terminale]
$ tsc constraints.ts
constraints.ts:1:16 - error TS2564: Property 'species' has no initializer and is not definitely assigned in the constructor.

1 class Animal { species: string; }
                 ~~~~~~~

constraints.ts:2:30 - error TS2564: Property 'canFly' has no initializer and is not definitely assigned in the constructor.

2 class Bird extends Animal  { canFly: boolean; }
                               ~~~~~~

constraints.ts:3:30 - error TS2564: Property 'isVenomous' has no initializer and is not definitely assigned in the constructor.

3 class Snake extends Animal { isVenomous: boolean; }
                               ~~~~~~~~~~


Found 3 errors in the same file, starting at: constraints.ts:1

\end{console}

Basterebbe dare un valore iniziale alle proprietà, come in
\code{canFly = false}, oppure dichiararle con un punto esclamativo, come in
\code{species!: string}: vuol dire «verrà assegnata prima dell'uso, me ne
occupo io». Senza controlli severi il file compila come nelle slide:

\begin{console}[Terminale]
$ tsc --strict false constraints.ts
$ node constraints.js
Dog
Kiwi
Cobra
Spider
\end{console}

\section{Altri tipi da conoscere}

TypeScript include altri tre tipi, utili soprattutto con le funzioni.

\begin{center}
\renewcommand{\arraystretch}{1.45}
\begin{tabular}{@{}l@{\hspace{1.2cm}}>{\raggedright\arraybackslash}p{8.6cm}@{}}
  \textbf{Tipo} & \textbf{Significato} \\
  \code{void}    & tipo di ritorno delle funzioni che non restituiscono un
                   valore \\
  \code{unknown} & rappresenta qualunque valore possibile: simile ad
                   \code{any}, ma \textbf{più sicuro}, perché non consente
                   alcuna operazione senza prima restringere il tipo \\
  \code{never}   & rappresenta valori che non vengono mai osservati: il tipo di
                   una funzione che non ritorna mai, o di un ramo irraggiungibile \\
\end{tabular}
\end{center}

\code{void} è anche il tipo dedotto per le funzioni senza \code{return}, o con
un \code{return} vuoto:

\begin{lstlisting}[style=ts, name=void.ts]
function f(): void {
  return;
}

function g() {            // function g(): void
  console.log("Hello");
}

let x: void = g();
\end{lstlisting}

\begin{lstlisting}[style=ts, name=unknown.ts]
function greet(a: any) {
  console.log(a.name);       // compila, ma possibile errore a runtime
}

function saferGreet(a: unknown) {
  console.log(a.name);       // errore di compilazione
}

type namedObject = { name: string };

function notSoSafeGreet(a: unknown) {
  console.log((a as namedObject).name);   // asserzione esplicita
}
\end{lstlisting}

\begin{console}[Terminale]
$ tsc unknown.ts
unknown.ts:6:15 - error TS18046: 'a' is of type 'unknown'.

6   console.log(a.name);       // errore di compilazione
                ~


Found 1 error in unknown.ts:6

$ tsc --strict false unknown.ts
unknown.ts:6:17 - error TS2339: Property 'name' does not exist on type 'unknown'.

6   console.log(a.name);       // errore di compilazione
                  ~~~~


Found 1 error in unknown.ts:6

\end{console}

L'errore \code{TS18046} delle slide è quello della modalità severa; senza,
TypeScript dice la stessa cosa in un altro modo: la proprietà \code{name} non
esiste sul tipo \code{unknown} (\code{TS2339}). In entrambi i casi per usare il
valore bisogna prima restringerne il tipo, o ricorrere a un'asserzione come in
\code{notSoSafeGreet}, che però non controlla nulla.

\begin{lstlisting}[style=ts, name=never.ts]
function f(x: string | number) {
  if (typeof x === "string") {
    return x.toLowerCase();
  } else if (typeof x === "number") {
    return x.toPrecision(2);
  } else {
    return x;   // qui x ha tipo never: nessun altro caso e' possibile
  }
}

function g(x: any): never {   // g non ritorna mai: solleva sempre un errore
  throw new Error("Oops!");
}
\end{lstlisting}

\section{Livelli di severità}

Utenti diversi si aspettano cose diverse da TypeScript. Nelle slide, cioè
fino a TypeScript 5, l'esperienza predefinita è \textbf{opzionale}: i tipi non
sono obbligatori e si ricorre ad \code{any} quando non è possibile dedurre un
tipo preciso. Il flag \code{noImplicitAny} rende più severo questo
comportamento.

Per non essere troppo invasivo, TypeScript 5 tiene inoltre spenti di partenza
alcuni controlli. Per esempio \code{null} e \code{undefined} possono essere
assegnati a qualunque tipo: dimenticarsi di gestirli esplicitamente è la causa
di innumerevoli bug, tanto che il concetto è stato definito \guillemotleft
l'errore da un miliardo di dollari\guillemotright. Il flag
\code{strictNullChecks} impone che \code{null} e \code{undefined} siano gestiti
in modo esplicito.

\begin{lstlisting}[style=ts, name=strict.ts]
class CustomCollection<T> {
  list: T[] = [];
  add: (x: T) => number = (element: T) => {
    return this.list.push(element);
  }
}

let x = new CustomCollection<string>();
x.add("Sam"); x.add("Cliff"); x.add("Mama");

let character = x.list.find((val: string) => {
  return val === "Amelie"
});
// character potrebbe essere undefined!
console.log(character.toUpperCase());
\end{lstlisting}

\code{find} restituisce il primo elemento che soddisfa la condizione, oppure
\code{undefined} se non ne trova: è il caso di \code{"Amelie"}. Ecco i comandi
delle slide, eseguiti con TypeScript 5:

\begin{console}[Terminale]
$ npx --yes -p typescript@5.9.3 tsc strict.ts
strict.ts:11:24 - error TS2550: Property 'find' does not exist on type 'string[]'. Do you need to change your target library? Try changing the 'lib' compiler option to 'es2015' or later.

11 let character = x.list.find((val: string) => {
                          ~~~~


Found 1 error in strict.ts:11

$ npx --yes -p typescript@5.9.3 tsc -target es6 strict.ts
$ npx --yes -p typescript@5.9.3 tsc -target es6 -strictNullChecks strict.ts
strict.ts:15:13 - error TS18048: 'character' is possibly 'undefined'.

15 console.log(character.toUpperCase());
               ~~~~~~~~~


Found 1 error in strict.ts:15

\end{console}

Il primo comando spiega perché nelle slide compare \code{-target es6} (le
opzioni si possono scrivere anche con un trattino solo): \code{find} è
arrivato con ES6, e con il bersaglio predefinito di TypeScript 5, ES5, il
compilatore non lo conosce. Il secondo compila senza errori; il terzo, con
\code{strictNullChecks}, intercetta il bug. (Nelle slide l'errore è alla riga
13: il file compilato lì era evidentemente un po' più corto di quello
mostrato.)

Ecco che cosa succede se il bug arriva a runtime:

\begin{console}[Terminale]
$ tsc --strict false strict.ts
$ node strict.js
[cartella del progetto]/strict.js:16
console.log(character.toUpperCase());
                      ^

TypeError: Cannot read properties of undefined (reading 'toUpperCase')
    at Object.<anonymous> ([cartella del progetto]/strict.js:16:23)
    ...

Node.js v22.22.0
\end{console}

Da TypeScript 6.0 l'impostazione predefinita è cambiata: l'opzione
\code{strict} è attiva di partenza, e con lei tutti i controlli della famiglia,
fra cui \code{noImplicitAny}, \code{strictNullChecks} e
\code{strictPropertyInitialization} (quello delle proprietà senza valore
iniziale). Con TypeScript 7 il bug viene quindi intercettato anche senza flag:

\begin{console}[Terminale]
$ tsc --strict false --strictNullChecks strict.ts
strict.ts:15:13 - error TS18048: 'character' is possibly 'undefined'.

15 console.log(character.toUpperCase());
               ~~~~~~~~~


Found 1 error in strict.ts:15

$ tsc strict.ts
strict.ts:15:13 - error TS18048: 'character' is possibly 'undefined'.

15 console.log(character.toUpperCase());
               ~~~~~~~~~


Found 1 error in strict.ts:15

\end{console}

\begin{center}
\renewcommand{\arraystretch}{1.45}
\begin{tabular}{@{}>{\raggedright\arraybackslash}p{6.2cm} l l@{}}
  \textbf{Impostazione predefinita} & \textbf{TypeScript 5} & \textbf{TypeScript 6 e 7} \\
  bersaglio (\code{target}) & ES5 (ES3 fino alla 4.9) & ES2025 \\
  modalità severa (\code{strict}), con \code{noImplicitAny},
  \code{strictNullChecks} e gli altri controlli & disattiva & attiva \\
\end{tabular}
\end{center}

Senza il controllo il codice compila e il bug esplode a runtime; con il
controllo viene intercettato subito. È il motivo per cui nei progetti nuovi
conviene lavorare fin dall'inizio in modalità severa: dalla versione 6 è
TypeScript stesso a sceglierla per noi.

\section{Riferimenti}

\begin{itemize}
  \item \textbf{The TypeScript Handbook}. \\
        \urlx{https://www.typescriptlang.org/docs/handbook/intro.html} \\
        Non serve conoscerlo per intero: va usato come riferimento per le parti
        discusse in questo capitolo.
  \item \textbf{Lecture Notes for the Comparative Studies of Programming
        Languages Course}, Paquet e Mokhov (2010). \\
        \urlx{https://arxiv.org/pdf/1007.2123.pdf} \\
        Sezione 3.3 per approfondire i sistemi di tipi (duck typing, sistemi
        strutturali contro nominativi). Non richiesto per il corso.
  \item \textbf{To Type or Not to Type? A Systematic Comparison of the Software
        Quality of JavaScript and TypeScript Applications on GitHub}, Justus
        Bogner e Manuel Merkel (2022). \\
        \urlx{https://dl.acm.org/doi/pdf/10.1145/3524842.3528454} \\
        Lettura interessante: gli autori confrontano applicazioni JavaScript e
        TypeScript sotto molti aspetti, dalla qualità e leggibilità del codice
        alla propensione ai bug e ai tempi di correzione.
\end{itemize}
